% ===== §6 When Ethics Reaches Bottom, It Finds Aesthetics =====
\section{The Genealogy of the Aesthetic Grounding of Ethics}
\label{p22:sec:genealogy}

\epigraph{\zh{兴于诗，立于礼，成于乐。}}{\textit{Analects}}

\noindent
The reduction chain reaches, at its third link, a claim that is easy to state and easy to mistake for a stipulation. The construction of a field, the chain says, is an aesthetic matter; the aesthetic and the ethical are coupled at their root; therefore the cultivation of the field-building capacity is an aesthetic education that is at the same time an ethical one. The danger is that this coupling, the joint of aesthetic and ethical on which the whole chain turns, will look like something the paper has simply defined into being, by christening a single sensibility that judges both the appropriate and the good and then announcing that the two are therefore one. A coupling asserted by definition is worth nothing, since one may define as one pleases and the world is not obliged to agree. This section exists to show that the coupling is not the paper's invention. It is a recognition that a long line of thinkers, working in different centuries, idioms, and traditions, have arrived at independently and from opposite directions, the recognition that when ethics is pressed to its ground it discovers aesthetics beneath it, that value is in the first place perceived rather than inferred, and that the appropriate overflows every rule that would contain it. The paper's own contribution, set out after the witnesses, is not the coupling but a mechanism for it; the coupling itself is an inheritance, and the genealogy is the proof of inheritance.

\subsection{Confucius: the person is completed in the aesthetic}
\label{p22:sub:gen-confucius}

The oldest witness is also the most direct, and it belongs to the roots this series has drawn on throughout. The formation of a person, the \textit{Analects} says, is begun in poetry, established in rite, and completed in music~\parencite{confucius_analects}: the sequence does not end in knowledge or in action but in an aesthetic accomplishment, and the completion is not decorative but constitutive, the final forming of the ethical person in and as a refinement of feeling. The tradition that follows does not treat this as a metaphor. Music, in the rite-and-music teaching, is the means by which the dispositions are tuned, by which a person and a people are brought into accord, so that the cultivation of feeling and the cultivation of virtue are one undertaking carried out in one medium. What is striking, set against the later Western quarrels the section will reach, is that this tradition never had to discover the coupling of the aesthetic and the ethical, because it never divided them; the good and the beautiful were given together, and the labour was not to join them but to keep them joined. The Confucian witness therefore stands at the head of the genealogy not as the first to argue the coupling but as the tradition for which the coupling was never in question, the place from which the later discoveries look like the recovery of something that had been forgotten.

\subsection{Schiller: the moral problem is first an aesthetic one}
\label{p22:sub:gen-schiller}

The founding modern witness made the coupling into a programme. Confronting the political and moral wreckage of his age, he argued that the way to the moral condition does not run directly through the moral, that one cannot make people good by commanding goodness or by perfecting their reason, because a humanity split between sensuous compulsion and rational compulsion is incapable of the freedom that morality requires~\parencite{schiller_aesthetic}. The route to the moral runs through the aesthetic, because it is beauty, and the play it awakens, that releases the person from both compulsions at once and restores the wholeness in which moral agency becomes possible. The moral problem, on this account, must first be solved as an aesthetic problem, which is precisely the direction of reduction this paper follows. The witness is founding because it is the first to say, as a thesis and not an aside, that aesthetic education is not one province of education among others but the condition of the ethical, the cultivation without which the ethical cannot be reached. Where this paper parts from it is marked and not hidden: the founding account locates the reconciliation within the individual, the play that unifies one person's faculties, whereas the field-building this paper describes is relational from the start, the aesthetic that grounds the ethical between subjects and not within one. The genealogy takes the thesis and relocates its site, from the single psyche to the coupling.

\subsection{Wittgenstein: ethics and aesthetics are one}
\label{p22:sub:gen-wittgenstein}

The sharpest witness is a single line, and its sharpness is worth the section's care. Ethics and aesthetics are one, the early work says~\parencite{wittgenstein_tractatus}, and the unity is not a slogan but a consequence of the work's deepest commitment, that value does not lie within the world among the facts and so cannot be stated in propositions that picture facts. Both the ethical and the aesthetic belong to what cannot be said but only shown, to the domain of how the world is regarded rather than what is the case within it, and they are one because they are the same standing-outside-the-facts, the same way of looking upon the world as a whole rather than describing its contents. This is the most exact possible statement of the joint this paper needs, and from an unexpected quarter, since it arrives not from a philosophy of feeling but from the most austere logic. What it gives the genealogy is the claim that appropriateness, the thing the field-builder must perceive, cannot be captured in propositions and rules precisely because it shares its nature with value as such, which is shown and not said. The line is the genealogy's keenest blade because it makes the coupling not a psychological observation about how the good and the beautiful happen to feel alike, but a structural claim about where both lie, outside the sayable, in the showing.

\subsection{Scheler: value is felt before it is reasoned}
\label{p22:sub:gen-scheler}

If the previous witness placed value outside the sayable, the next says how it is nonetheless apprehended, and the answer feeds directly into this paper's account of perception. Moral value, on this witness's account, is not concluded from premises but felt, grasped in a feeling-of-value that is a mode of perception in its own right, prior to and independent of intellectual inference; there is an ordering of values given to the heart that the understanding does not establish but discovers~\parencite{scheler_formalism}. This is the philosophical charter for the claim, central to this paper's coupling, that the sensibility which perceives appropriateness is the same one that perceives the good, because on this account the good is in the first place something perceived, a value felt in an act of value-perception rather than derived in an act of reasoning. The witness matters most to the section on perception, where the seeing of the field's appropriateness is described, because it establishes, independently of this paper, that value-apprehension is perceptual at its root, that to know the good is first to feel it as one feels a quality, and so that a sensibility trained to perceive appropriateness is not metaphorically but literally a moral organ. The coupling, on this witness, is grounded in the structure of value-apprehension itself: the good is felt, and what feels it is a perception, and a perception refined upon the appropriate is thereby refined upon the good.

\subsection{Gadamer: taste was a moral power before it was an aesthetic one}
\label{p22:sub:gen-gadamer}

The next witness recovers a history that the section's own key term has half-forgotten. Before taste was narrowed into a category of art appreciation, it named a power of the common sense, a capacity of judgement that found the right and fitting in matters where no rule could be applied, a moral and social discernment of what is appropriate to a situation~\parencite{gadamer_truth}. Judgement of this kind, the capacity to see what fits where no rule determines the fitting, was understood as continuous with moral judgement and not separate from it, and the narrowing by which taste became a purely aesthetic affair, severed from the discernment of the appropriate in conduct, is diagnosed by this witness as a loss, a forgetting of taste's original breadth. The relevance to this paper is exact and double. The recovered concept of taste, the discernment of the fitting where no rule applies, is almost word for word the appropriateness this paper says the field-builder must perceive and the engineering account cannot capture; and the diagnosis of taste's narrowing tells this paper that to return aesthetic education to ethics is not an innovation but a restoration, the undoing of a forgetting. The witness thus gives the genealogy its sense of recovery: the coupling the paper asserts was once common knowledge, lodged in the very word taste, and was lost when aesthetics was sealed off as a domain of its own.

\subsection{Williams and Murdoch: the analytic recovery of perception over theory}
\label{p22:sub:gen-williams-murdoch}

The genealogy's final witnesses come from within the analytic tradition itself, which matters because they show the recognition is not the property of one philosophical school but is reached, independently, even by those least disposed to mysticism about value. The first of them mounted a sustained case against the ambition of moral theory, the project of capturing ethical life in principles and decision procedures, arguing that this ambition mistakes the nature of ethical life, which depends on perceptions and judgements in the concrete situation that no theory can anticipate or replace, so that the very drive to a moral theory is a kind of error~\parencite{williams_limits}. This is direct analytic support for the link in this paper's chain that says appropriateness cannot be reduced to rules, and it arrives from the tradition most committed to rules and arguments, which is why it carries the weight it does. The second of these witnesses gave the recognition its most beautiful form, making attention the centre of the moral life: to act well is first to see well, to attend to the other and the situation justly and without the self's distorting interference, and this just attention is described in frankly aesthetic terms, modelled on the selfless looking that art demands~\parencite{murdoch_sovereignty}, so that the moral gaze and the aesthetic gaze are shown to be one kind of seeing. This last witness is the most direct ancestor of this paper's claim that the perception of the field is at once aesthetic and moral, since it says, as plainly as the paper could wish, that to see justly is the moral act and that such seeing is of a piece with the contemplation of beauty. With these two the genealogy reaches the present, having shown the coupling recognised across the whole breadth of philosophy, from the oldest Eastern teaching to the most recent analytic ethics, by thinkers who shared almost nothing except this.

\subsection{The dissenting witness and the danger it marks}
\label{p22:sub:gen-dissent}

The genealogy would be dishonest if it gathered only the witnesses who support the coupling and suppressed the one who turns it against the paper's purpose, and the suppressed witness is instructive enough to be admitted. There is a line of thought that affirms the coupling of the aesthetic and the ethical and draws from it exactly the opposite of what this paper draws, that takes the priority of the aesthetic to mean the subordination of the ethical to it, the justification of life as an aesthetic phenomenon in which the demands of justice are dissolved into the demands of beauty, the strong and the beautiful released from a morality recast as the resentment of the weak. This witness agrees that ethics finds aesthetics at its bottom and concludes that ethics is therefore to be overcome in the name of the aesthetic, that the fine life is beyond the categories of the good and the just, that beauty is the highest value and answers to nothing above it. The genealogy must admit this witness because it shares the premise the genealogy draws on, the priority of the aesthetic, while drawing from it a conclusion the paper must refuse, and the refusal is not a matter of preferring the agreeable witnesses but of showing that the dissenting conclusion does not follow from the shared premise.

It does not follow, and the orthogonality the next section establishes is exactly why. The dissenting witness moves from the priority of the aesthetic to the supremacy of the aesthetic, from the claim that value is perceived before it is reasoned to the claim that the perception of beauty is the whole of value and the just a mere constraint to be transcended. But the priority of the aesthetic, as this paper understands it, is the priority of one perception, the perception of the field's generative direction, and that perception, the orthogonality will show, is silent on justice, on who is rendered fuel for the accumulation it perceives. To make the aesthetic supreme, to release the beautiful life from the demands of justice, is precisely to commit the error the orthogonality guards against, the elevation of the perception of accumulation over the question of who pays for it, the aestheticism that the exquisite cruelty exemplifies. The dissenting witness is therefore not a counter-authority that the genealogy must defeat but an illustration of the danger the whole paper is built to meet, the danger that the real coupling of the aesthetic and the ethical, misunderstood as the supremacy of the aesthetic, licenses the fine and the unjust. The genealogy includes the dissenter to mark this danger at its source, and to show that the paper's coupling, bounded by the orthogonality, yields the opposite of the dissenter's conclusion: not the transcendence of justice by beauty but the insistence that beauty, even the deepest relational beauty, is necessary and not sufficient for the good, and that the perception which finds the appropriate must be joined by the separate regard that asks who is consumed.

\subsection{The disagreements among the witnesses}
\label{p22:sub:gen-disagreement}

Before the convergence is claimed it must be made honest, because the witnesses disagree about almost everything except the one thing the genealogy draws from them, and a convergence asserted over suppressed disagreement would be worth no more than a stipulation. They disagree about the nature of value: for one it is felt in an objective order given to the heart, for another it is shown and cannot be said at all, for another it is recovered as a social judgement sedimented in a tradition, for another it is the selfless attention that art teaches. They disagree about method: the austere logician and the phenomenologist of value and the analytic critic of theory would each find the others' procedures suspect, and the Confucian completion of the person in music belongs to a form of life the modern witnesses do not share and might not recognise. They disagree about whether there is a moral knowledge at all, about whether value is in the world or in the regard brought to it, about nearly every question a philosophy of value can raise. The genealogy does not paper over this. It claims convergence on one point only, and the value of the convergence depends on its being convergence from genuinely different places, since a point reached independently by thinkers who agree on nothing else is far better evidenced than a point on which a single school insists. The disagreement is not the genealogy's embarrassment but its strength: that so many, sharing no method and no metaphysics, were nonetheless driven to the same recognition about value and rules is precisely what makes the recognition hard to dismiss as the artefact of one school's commitments.

This is the method of the series applied to the genealogy itself, the seeking of common ground while the differences are kept. The differences among the witnesses are kept, named rather than dissolved, and the common ground is sought within them, the one recognition they share across everything that divides them. To force them into a single doctrine would be to reduce the many to one in exactly the way the paper has refused throughout; to leave them merely juxtaposed, a list of opinions about value, would be to miss the convergence that is the genealogy's whole point. The genealogy holds the two together, the kept differences and the sought common ground, and what it finds in the common ground is set out next.

\subsection{The limit of each witness, and the critiques they have drawn}
\label{p22:sub:gen-limits}

The genealogy gathers the witnesses for the one recognition they share, but it would flatter them to present each as simply correct, and the dialectical-materialist commitment that governs this paper requires that each be given its limit and that the critiques each has drawn be registered, since a recognition is better held when one knows the failures of those who reached it. The witnesses are connected not only by their shared recognition but by a shared liability, the liability of taking the priority of perceived value to license an inattention to the material and political conditions under which perception is formed, and this liability takes a specific form in each.

The Confucian completion of the person in music has drawn the critique that its harmony is the harmony of a hierarchical order, that the rite-and-music tradition tuned the dispositions to a settled social arrangement and called the tuning virtue, so that its seamless union of the beautiful and the good can sanctify a domination as easily as a flourishing. The critique is just to this extent, that the tradition did not by itself supply the orthogonality that distinguishes a generative concord from a beautiful subjection, and the paper supplies what the tradition lacked, the question of who is rendered fuel for the harmony, which the rite-and-music tradition did not systematically raise. Schiller has drawn the critique of an aestheticism avant la lettre, that his aesthetic state, by reconciling the faculties in play, risks a withdrawal from the political struggle whose conditions it claims to prepare, a beautiful soul cultivated while the unjust order stands; and the critique bites against any reading of aesthetic education, including this paper's, that would let the cultivation of the field substitute for the change of the conditions, which is why the political economy of the field was necessary. Wittgenstein's identification of ethics and aesthetics, placing both outside the sayable, has drawn the critique that it removes value from rational discussion altogether, consigning the good and the beautiful to a showing that cannot be argued, and the paper parts from him exactly here, supplying the mechanism that brings the coupling back within reasoned articulation, since the holonomy of a field is a structure about which one can reason, not only a value that can be shown.

Scheler's intuition of a given order of values has drawn the critique of a material aprioris that no two cultures share, the objection that the values felt as given differ across forms of life, so that the feeling-of-value cannot ground a universal ethics; and the paper accommodates the critique by grounding value not in a given order felt by intuition but in the generative direction of a real field, which is variable in its content across relations and constant in its form, the accumulation of holonomy. Gadamer's recovery of taste as moral judgement has drawn the critique that the tradition whose sensus communis he restores is itself a particular and exclusionary inheritance, that the common sense of the fitting is the common sense of a class and a culture, and the paper registers this in its political economy, where the conditions of the sensibility's cultivation are shown to be unequally distributed. Williams's case against moral theory has drawn the critique that, in refusing theory, it leaves ethics without the critical leverage that principles provide against a corrupt practice, that a perception trained within an unjust order will perceive its injustice as fitting; and the paper answers with the orthogonality, the deliberate question of justice that the trained perception does not contain and must be made to ask. Murdoch's attention has drawn the critique that just attention, however pure, can attend lovingly to a scene whose injustice it does not register, that the selfless gaze is not by itself a just one; and this is the orthogonality again, met now as the limit of even the most morally serious of the witnesses, that fineness of attention is not justice of direction.

What this accounting shows is that the witnesses converge not only on a recognition but on a shared limit, the limit at which the priority of perceived value, unguarded, becomes an inattention to the conditions and the justice that perception alone does not deliver. The paper inherits the recognition and corrects the limit, supplying at the level of mechanism what the witnesses left at the level of intuition, and supplying through the orthogonality and the political economy the attention to justice and condition that the witnesses, severally, neglected. This is the materialist completion of the genealogy: the recognition that ethics finds aesthetics at its bottom is preserved, and the idealist temptation that has shadowed every witness, the temptation to let the priority of the aesthetic excuse an inattention to material conditions and to justice, is refused.

\subsection{The common structure}
\label{p22:sub:gen-common}

The witnesses share almost nothing in method, idiom, or century, and the worth of the genealogy lies precisely in that they nonetheless converge. Drawn together, their convergence has two strands. The first is a claim about where value lies: that the good is not in the first place inferred from principles but perceived, felt, discerned, apprehended in an act more like seeing than like reasoning, whether named the feeling-of-value, the just attention, the taste that finds the fitting, or the music in which the person is completed. The second is a claim about rules: that the appropriate, the fitting, the right-in-this-situation, overflows every rule that would contain it, so that ethical life cannot be captured in a theory or a procedure and depends instead on a discernment exercised in the concrete, a discernment continuous with the discernment of beauty. These two strands are the two halves of one recognition. If value is perceived rather than inferred, then the faculty of value is a perceptual one; and if the appropriate overflows rules, then that perceptual faculty is needed exactly where rules give out, which is the domain the witnesses agree to call, in its widest sense, the aesthetic. When ethics is pressed to its ground, it finds beneath it not a deeper rule but a way of perceiving, and a way of perceiving the fitting where no rule decides is what aesthetic discernment has always been. That is the inheritance, and it is not one school's property but a recurrence across the whole field of reflection on value.

\subsection{The asymmetry of East and West in the genealogy}
\label{p22:sub:gen-asymmetry}

A feature of the genealogy deserves remark before its contribution is stated, the asymmetry between how the coupling appears in the Eastern and the Western witnesses, because the asymmetry is instructive about the coupling itself. In the Western witnesses the coupling appears as a discovery or a recovery, something arrived at against a prevailing separation, won back from a tradition that had divided the aesthetic from the ethical and needed to be shown their unity. Schiller argues for the coupling against a moralism that would reach the good directly; Wittgenstein states it as a hard-won consequence of a difficult logic; Gadamer recovers it against the narrowing that had severed taste from moral judgement; the analytic witnesses reach it against a theory-driven ethics that had forgotten perception. In each the coupling is the conclusion of a struggle against a separation, the reunification of what the Western tradition had put asunder. In the Eastern witnesses the coupling appears otherwise, not as a discovery or recovery but as an inheritance never lost, a unity that was never divided and so never needed to be won back. The Confucian completion of the person in music, the seamless joining of the beautiful and the good in the rite-and-music tradition, present the coupling as simply the way things are, never having been separated, never requiring an argument to reunite them.

This asymmetry is not a defect of the genealogy but a clue to the coupling's status, and reading it correctly strengthens rather than weakens the case. That the Western tradition had to win the coupling back, against its own separation of the aesthetic from the ethical, shows the separation to be a contingent achievement of that tradition and not a feature of the things themselves; the West divided what need not have been divided, and its best thinkers spent their effort recovering the unity its division had obscured. That the Eastern tradition never divided them shows the unity to be available without the division, a way of holding the aesthetic and the ethical together that some traditions simply kept. The two together make the strongest possible case for the coupling: the tradition that divided them found, at its most rigorous, that it had to put them back together, and the tradition that never divided them shows that the division was never necessary. The coupling is thus confirmed both by those who recovered it against their separation and by those who never lost it, and the asymmetry between recovery and inheritance, far from weakening the genealogy, doubles its force, since a unity reached by recovery in one tradition and held as inheritance in another is better evidenced than one attested in only a single way. The paper's own contribution, the mechanism, is offered to both: to the West it gives the structure that explains why the division it recovered from was always artificial, and to the East it gives the formal articulation of a unity it held without needing to formalize, the holonomy and the field that say, in the language of mechanism, what the rite-and-music tradition knew without saying.

\subsection{The gift and the reservation of the genealogy}
\label{p22:sub:gen-contribution}

The genealogy gives this paper its licence and withholds its mechanism, and the distinction between the two is what the paper's own contribution turns on. What the witnesses establish is that the coupling of the aesthetic and the ethical is real, recurrent, and independently arrived at, so that to assert it is to recover a recognition rather than to stipulate a definition. This is the licence, and it is what the reduction chain's third link needed, since a link that merely defined its coupling would have carried no weight, whereas a link that recovers a coupling recognised from Confucius to the present carries the weight of the whole tradition behind it. What the witnesses do not give is an account of how the coupling works, by what mechanism the perception of appropriateness and the perception of the good are one perception, in what medium a sensibility refined upon the fitting becomes thereby a moral organ. They report the coupling; they do not explain it. Each names the fact from his own side, the feeling-of-value, the just attention, the completing music, but none supplies the structure that would say why value is perceptual and why the appropriate overflows rules, why these two strands belong together rather than merely co-occurring.

\claim{The coupling of the aesthetic and the ethical is an inheritance, recognised independently across traditions from the Confucian completion of the person in music to the analytic recovery of perception over theory; it is not this paper's stipulation. What the tradition reports but does not explain is the mechanism of the coupling. The contribution of this paper is to supply that mechanism: appropriateness is the perceptible signature of accumulating holonomy in a relational field, and value is that accumulation, so that to perceive the appropriate and to perceive the good are one perception because the appropriate is what the good looks like to perception. The genealogy supplies the licence; the framework supplies the mechanism; neither alone would suffice.}

The next section sets out that mechanism, reading the coupled sensibility as the resolution of a relational field and the perception of appropriateness as the registration of its accumulating holonomy, so that what the witnesses reported from the outside is given an inside. One guard must be carried forward into that account, and it is the guard the following section is built around. To say that the perception of the appropriate and the perception of the good are one perception is not to say that refinement of the sensibility guarantees the good, for a sensibility can be exquisitely refined and exquisitely unjust, and the genealogy's witnesses, who attended to the fineness of perception, did not all attend to its direction. The coupling of the aesthetic and the ethical is a coupling of one perception to one register of value, not a guarantee that the perceptive are the just; and the orthogonality of refinement and justice, which the next section but one must establish, is what keeps this genealogy from becoming the aestheticism it might otherwise be mistaken for. The recognition that ethics finds aesthetics at its bottom is the beginning of the account and not its end, because what is found there must still be shown to have a direction, and a fineness without a direction is precisely the danger the whole of this paper is built to meet.
